\section{Тестирование программного комплекса}

\subsection{Стратегия тестирования}

Тестирование программного комплекса проводилось на трёх уровнях:
\textit{модульном} (проверка корректности отдельных функций),
\textit{интеграционном} (проверка взаимодействия модулей через файловый интерфейс)
и \textit{системном} (проверка сквозного выполнения полного пайплайна).

\subsection{Модульные тесты}

\subsubsection{Тестирование вычисления реализованной волатильности}

Для оценщиков Garman--Klass и Parkinson проверялись:
\begin{enumerate}
    \item \textbf{Граничные случаи.}
          При $H_t = L_t = O_t = C_t$ (свеча без движения)
          оба оценщика возвращают $\widehat{\sigma}^2_t = 0$.
    \item \textbf{Масштабирование.}
          При умножении всех цен на константу $k > 0$
          логарифмическая волатильность не изменяется:
          $\widehat{\sigma}^2(kO, kH, kL, kC) = \widehat{\sigma}^2(O, H, L, C)$.
    \item \textbf{Воспроизводимость.}
          Повторный запуск \texttt{compute\_rv.py} на тех же данных
          даёт побайтово идентичный CSV-файл.
    \item \textbf{Сравнение оценщиков.}
          На синтетических данных, сгенерированных из геометрического
          броуновского движения с известной дисперсией $\sigma^2 = 0{,}04$,
          среднее значение оценщика Garman--Klass отклоняется от истинного
          не более чем на $5\%$ при $n > 250$ наблюдений.
\end{enumerate}

\subsubsection{Тестирование парсеров}

Каждый парсер тестировался следующим образом:

\begin{itemize}
    \item \textbf{download\_moex.py:} запрос на однодневный период возвращает
          ровно одну строку OHLCV; колонки \texttt{open, high, low, close, volume}
          присутствуют и имеют корректные типы данных (\texttt{float64}).
    \item \textbf{parse\_cbr.py:} курс USD/RUB на 01.01.2022 соответствует
          официальному значению ЦБ РФ (74{,}29 руб.) с точностью до копейки.
    \item \textbf{parse\_fred.py:} значение VIX на 24.02.2022 превышает 30{,}0
          (проверка корректности загрузки данных о волатильности в период шока).
    \item \textbf{aggregate\_sentiment.py:} результирующий датафрейм содержит
          ровно 25 признаков; доля NaN-значений равна 0 после forward fill.
\end{itemize}

\subsubsection{Тестирование HAR-модели}

\begin{enumerate}
    \item \textbf{Детерминизм.}
          Модель HAR-log при двух последовательных запусках
          на одних данных возвращает идентичные коэффициенты.
    \item \textbf{Структура предсказания.}
          При expanding window с начальным окном 120 дней
          длина вектора прогнозов совпадает с числом тестовых наблюдений.
    \item \textbf{Знаки коэффициентов.}
          На реальных данных GLDRUB\_TOM все три коэффициента
          $\hat\beta_d, \hat\beta_w, \hat\beta_m > 0$,
          что соответствует теоретическим ожиданиям модели HAR-RV.
\end{enumerate}

\subsubsection{Тестирование XGBoost}

\begin{enumerate}
    \item \textbf{Воспроизводимость.}
          Фиксация \texttt{random\_state=42} обеспечивает идентичные метрики
          при повторных запусках.
    \item \textbf{Отсутствие утечки данных.}
          Разбиение 60/40 производится строго по времени:
          ни один тестовый индекс не присутствует в обучающей выборке.
    \item \textbf{Корректность метрик.}
          Значение $R^2_{\mathrm{OOS}}$ вычисляется относительно
          исторического среднего обучающей выборки,
          а не тестовой (проверка реализации формулы~Goyal--Welch).
\end{enumerate}

\subsection{Интеграционное тестирование}

Интеграционный тест запускает сквозной пайплайн на коротком
подпериоде данных (2023-01-01~--- 2023-12-31):

\begin{enumerate}
    \item \textbf{download\_moex.py} → \texttt{GLDRUB\_TOM\_1d.csv} (248 строк)
    \item \textbf{compute\_rv.py} → \texttt{rv\_features.csv} (246 строк, 10 признаков)
    \item \textbf{aggregate\_sentiment.py} → \texttt{daily\_sentiment.csv} (365 строк, 25 признаков)
    \item \textbf{ml\_models.py} → \texttt{GLDRUB\_TOM\_xgb\_predictions.csv}
\end{enumerate}

Тест проходит успешно, если:
\begin{itemize}
    \item ни один из шагов не завершился с исключением;
    \item все выходные файлы созданы с ненулевым числом строк;
    \item $R^2_{\mathrm{OOS}}$ XGBoost на 2023~г.\ превышает $0{,}0$
          (модель превосходит наивный прогноз историческим средним).
\end{itemize}

\subsection{Системное тестирование}

Полный цикл обновления данных и переобучения моделей тестировался
на периоде 2018--2026 (производительность):

\begin{itemize}
    \item Загрузка котировок MOEX (4 инструмента): $\approx$\,45 сек.
    \item Загрузка CBR + FRED: $\approx$\,30 сек.
    \item Загрузка GDELT (инкрементальный режим, последние 30 дней): $\approx$\,5 мин.
    \item Агрегация sentiment: $<$\,10 сек.
    \item Обучение XGBoost (expanding window, 2 инструмента): $\approx$\,4 мин.
    \item Запуск Flask-сервера: $<$\,5 сек.
\end{itemize}

Суммарное время полного цикла составило $\approx$\,12 минут,
что удовлетворяет нефункциональному требованию
(не более 30 минут, раздел~\ref{sec:requirements}).

\subsection{Объёмные характеристики кода и данных}
\label{sec:volume}

Согласно требованиям к экспериментальной части программного проекта,
ниже приведены объёмные характеристики разработанного программного
комплекса и используемых данных
(таблицы~\ref{tab:vol_code} и~\ref{tab:vol_data}).

\begin{table}[h!]
\centering
\caption{Объём программного кода по подсистемам}
\label{tab:vol_code}
\begin{tabular}{lr}
\hline
Подсистема (директория) & Строк кода \\
\hline
\texttt{src/data/} (7 парсеров + агрегатор)        & 2\,378 \\
\texttt{src/models/} (HAR, XGBoost, kNN, RV, метрики) & 1\,440 \\
\texttt{src/nlp/}    (обработка текстов)            &    432 \\
\texttt{src/web/}    (Flask-роуты, шаблоны)         &    451 \\
\texttt{app.py}      (точка входа дашборда)         &    516 \\
\hline
\textbf{Всего Python-кода}                          & \textbf{5\,217} \\
\hline
\end{tabular}

\medskip
\raggedright
\footnotesize Суммарный объём исходного кода~--- $\approx$\,309 KB.
\end{table}

\begin{table}[h!]
\centering
\caption{Объём входных и промежуточных датасетов}
\label{tab:vol_data}
\small
\begin{tabular}{lrl}
\hline
Файл & Размер & Содержание \\
\hline
\texttt{data/processed/metals\_1d\_panel.csv}   & 310 KB & OHLCV 4-х инструментов \\
\texttt{data/processed/rv\_features.csv}         & 657 KB & RV + HAR-лаги (4\,238 строк) \\
\texttt{data/sentiment/daily\_sentiment.csv}     & 737 KB & 25 признаков $\times$ 3\,070 дней \\
\texttt{data/sentiment/cbr\_data.csv}            & 211 KB & Курсы и ставка ЦБ РФ \\
\texttt{data/sentiment/fred\_macro.csv}          & 336 KB & Brent, VIX, DXY, CPI, FED \\
\texttt{data/sentiment/google\_trends.csv}       & 246 KB & 3 группы запросов \\
\texttt{data/sentiment/gdelt\_sentiment.csv}     &  97 KB & Tone из GDELT GKG \\
\texttt{data/sentiment/rss\_news.csv}            &  66 KB & RSS 6 источников \\
\texttt{data/sentiment/telegram\_news.csv}       &  19 MB & 23\,733 сырых сообщений \\
\texttt{data/sentiment/telegram\_sentiment.csv}  &  59 KB & Дневная агрегация Telegram \\
\hline
\end{tabular}
\end{table}

Прогоны моделей в проекте сгенерировали 15 артефактов в
\texttt{data/processed/har\_results/} (прогнозы и метрики для трёх
спецификаций каждого из двух инструментов плюс файлы ablation)
и 6 артефактов в \texttt{data/processed/ml\_results/}
(XGBoost и kNN для GLDRUB/SLVRUB плюс сводные таблицы).
Сводная таблица попарного DM-теста (12 пар $\times$ 3~типа потерь)
сохраняется в \texttt{data/processed/dm\_test\_summary.csv}.

\subsection{Валидация данных}

Помимо функционального тестирования проводилась автоматическая
валидация качества данных:

\begin{enumerate}
    \item \textbf{Монотонность цен:}
          $H_t \geq \max(O_t, C_t)$ и $L_t \leq \min(O_t, C_t)$
          для каждой свечи. Нарушений не обнаружено.
    \item \textbf{Временной охват:}
          непрерывность дат в \texttt{daily\_sentiment.csv}
          (отсутствие пропущенных календарных дней) ~--- проверка пройдена.
    \item \textbf{Заполненность признаков:}
          доля NaN по каждому из 25 признаков равна 0~\% после агрегации.
    \item \textbf{Диапазон sentiment:}
          все значения \texttt{sentiment\_combined} $\in [-1{,}0;\, 1{,}0]$.
\end{enumerate}
